% Copyright 2026 Yoshida Shin
%
% This file is part of the ``Repeating Nines''.
%
% Permission is granted to copy, distribute, and/or modify this document
% under the terms of the GNU Free Documentation License, Version 1.3 or any later version published by the Free Software Foundation;
% with no Front-Cover Texts, no Back-Cover Texts, and one Invariant Section: ``This document was written by Shin Yoshida with the aim of contributing to a fairer and safer world.''
% A copy of the license is included in the section entitled ``GNU Free Documentation License.''


\begin{frame}[c]{}{}
  {\large Appendix}
\end{frame}

\begin{frame}[c]{}{}
  What did you think after hearing this explanation?

  \vspace{2ex}

  I believe that people who understand this kind of mathematical reasoning also have a different quality of thinking in daily work.
\end{frame}

\begin{frame}[c]{}{}
  For example, some of you might have been surprised that this explanation achieved its goal by building on a stack of obvious facts.
\end{frame}

\begin{frame}[c]{}{}
  I believe that this method of ``building on a stack of obvious facts'' is also useful in many fields.
\end{frame}

\begin{frame}[c]{}{}
  Also, when we considered the general shape of the graph for $a_n$, we thought about the case where n is extremely large.
\end{frame}

\begin{frame}[c]{}{}
  Thinking about extreme case to grasp the bigger picture is a powerful technique that I believe can be used in a wide range of fields.
\end{frame}

\begin{frame}[c]{}{}
  Sometimes, this way of thinking can be perceived as ``splitting hairs'' or ``being condescending'',

  \vspace{2ex}

  but I encourage you all to give it a try.
\end{frame}
